% Standalone mathematical section. Requires amsmath, amssymb, amsthm.
% Proofs independently audited against pinned sources; no computed premise
% occurs before the final subsection on the two-squared-axis certificate.
\section{Rigidity from dyadic fibers and prime-axis measurements}
\label{sec:primary-rigidity}

Throughout this section, a divisor ideal in a cyclic group means a union
of its subgroups. We use subgroup \emph{orders} when describing fibers.
Write the ambient group as
\[
G=\mathbb Z_{2^e}\times \mathbb Z_M,
\qquad M=\prod_{i=1}^n p_i^{a_i}\quad\hbox{odd}.
\]
The Chinese remainder theorem identifies this with the original labeled
cyclic group. At an odd point \(y\), the fiber of a divisor ideal \(A\)
is empty or a subgroup of \(\mathbb Z_{2^e}\): it is a union of members
of a chain. Let \(h_A(y)\in\{0\}\cup\{2^j:0\le j\le e\}\)
denote its size. If the odd
order of \(y\) divides that of \(z\), then \(h_A(y)\ge h_A(z)\).
Thus \(h_A\) is a nonincreasing function on the odd-order divisor lattice.
For a shell with order exponents \(v=(v_i)\), its population is
\(w(v)=\prod_i\varphi(p_i^{v_i})\).

On a \(p^a\)-axis, index a character by the order \(p^k\) of its
kernel. Its sum over the physical shell of order \(p^t\) is
\begin{equation}
R_p(k,t)=
\begin{cases}
\varphi(p^t),&t\le k,\\
-p^k,&t=k+1,\\
0,&t>k+1.
\end{cases}
\label{eq:primary-shell-row}
\end{equation}
For \(t=0\) the first line is one. This follows by subtracting the
character sums over the subgroups of orders \(p^t\) and \(p^{t-1}\).
Products of these rows compute all radial Fourier coefficients; in
particular these coefficients are integers. We use the unnormalized
Fourier transform throughout.

\paragraph{Character order versus kernel order.}
The character in~\eqref{eq:primary-shell-row} has conductor
\(p^{a-k}\), not \(p^k\). On a squarefree residual axis \(C_q\),
kernel exponent one means the trivial character and kernel exponent zero
means a nontrivial character. Thus a residual character of conductor
\(c>1\) has a zero kernel exponent on at least one axis. This convention
will be used explicitly in the proof of Theorem~\ref{thm:primary-power-four-primes}.

\begin{lemma}[Pure-axis measurements and common dyadic fibers]
\label{lem:common-dyadic-fibers}
Suppose that divisor ideals \(A,B\subseteq G\) have identical labeled
autocorrelations. Then the following hold.
\begin{enumerate}
\item Every Fourier coefficient supported on a single prime-power axis
is nonnegative, and it agrees for \(A\) and \(B\).
Consequently, for each fixed physical state on any odd axis, the total
signed population \(h_A-h_B\) over the complementary states is zero.
\item Wherever both dyadic fibers are nonempty, they are equal.
\item For every power of two \(L\), the number of odd points at which
the fibers are \((L,0)\) equals the number at which they are \((0,L)\).
Opposite signs of \(h_A-h_B\) cannot occur at comparable odd shells.
\end{enumerate}
\end{lemma}

\begin{proof}
Fix all coordinates except one prime-power coordinate. The resulting
section of a divisor ideal is empty or a subgroup of that cyclic
prime-power group. Its character sum is either zero or the order of the
subgroup. Summing over the fixed complementary coordinates proves
nonnegativity. Since homometry equates the squares of the real Fourier
coefficients, it equates these nonnegative coefficients themselves.
Fourier inversion on the chosen axis gives the stated marginal equality.

Fix an odd conductor \(c\). Let \(a_j(c)\) be the odd Fourier transform
of the indicator of \(\{y:h_A(y)=2^j\}\), and define \(b_j(c)\)
similarly. These are integers by \eqref{eq:primary-shell-row}. At a
dyadic character with kernel order \(2^k\), the two full Fourier
coefficients are
\[
P_k=\sum_{j=0}^k2^ja_j(c),\qquad
P'_k=\sum_{j=0}^k2^jb_j(c).
\]
Put \(D_k=P_k-P'_k\) and \(S_k=P_k+P'_k\), with
\(D_{-1}=S_{-1}=0\). Homometry says \(D_kS_k=0\) for every \(k\).
If \(a_j(c)-b_j(c)\) is odd, both increments from step \(j-1\) to
step \(j\) have exact 2-adic valuation \(j\). The incoming zero
component cannot be accompanied by another zero component, since then
both outgoing components would be nonzero. The incoming zero component
therefore becomes nonzero with valuation \(j\); it remains nonzero at
every later step, whose increments are divisible by \(2^{j+1}\).
The other component must remain zero. No later odd increment is
possible. Hence, at any fixed odd conductor, at most one size index
\(j\) has \(a_j(c)-b_j(c)\) odd.

If conclusion 2 failed, choose a minimal odd shell \(v\) with two
nonempty, unequal fibers. Every strictly smaller shell has two
nonempty equal fibers. On an axis with \(v_i>0\), choose kernel
exponent \(v_i-1\); when \(v_i=0\), choose the trivial character.
Modulo two, the product Fourier row is supported on the corners with
each coordinate either zero or \(v_i\). This follows from
\eqref{eq:primary-shell-row}, because \(p_i\) is odd and
\(\varphi(p_i^t)\) is even for \(t>0\). Every corner except \(v\)
has equal fibers. The two different nonzero sizes at \(v\) therefore
give two odd values of \(a_j(c)-b_j(c)\), contradicting the preceding
paragraph. This proves conclusion 2.

At the trivial odd character all \(P_k,P'_k\) are nonnegative, so
\(P_k=P'_k\). Subtract the equality at \(k-1\) from the equality
at \(k\); the counts of fibers of each exact size agree. Cancel the
common nonempty fibers using conclusion 2 to obtain conclusion 3.
Finally, if \(h_B(v)=0<h_A(v)\), then \(h_B\) vanishes at every
shell above \(v\), so an opposite discrepancy cannot occur there.
Interchanging \(A,B\) handles the other sign.
\end{proof}

\begin{corollary}[Two odd axes and a three-axis family]
\label{cor:two-odd-axes}
Every \(2^ep^aq^b\) is rigid, where \(p,q\) are distinct odd primes
and the exponents are nonnegative. Every \(2^ep^aq^br\) is rigid
when \(p,q,r\) are distinct odd primes and \(a,b\ge1\).
\end{corollary}

\begin{proof}
For two odd axes, fix an order state of the first axis. All states of
the second axis are comparable, so all discrepancies in this section
have the same sign by Lemma~\ref{lem:common-dyadic-fibers}. Their
weighted sum is zero by the pure-axis marginal equality, so there are
no discrepancies. The cases of fewer axes are included.

For three axes write their order states as \((i,j,k)\), where the
\(r\)-coordinate has \(k\in\{0,1\}\). Suppose a discrepancy exists
and take the least \(j\) occurring in any discrepancy. If there were
a discrepancy at \((i,j,0)\), marginal balance at the fixed state
\(i\) would require an opposite discrepancy \((i,j',k')\).
Comparability excludes \(k'=0\); with \(k'=1\), incomparability
requires \(j'<j\), contradicting the choice of \(j\). Thus all
discrepancies at this least \(j\) have \(k=1\). They form a chain
in \(i\), have one sign, and contradict marginal balance at \(j\).
\end{proof}

We call a shell \emph{disputed} if its two fibers differ. By the lemma,
exactly one fiber is then empty; its \emph{disputed size} is the size
of the other fiber. The following lemma is the main additional ingredient.

\begin{lemma}[Common zero along a changing primary axis]
\label{lem:primary-common-zero}
Let \(A,B\) be distinct homometric divisor ideals on
\(\mathbb Z_{2^e}\times\mathbb Z_{p^a}\times\mathbb Z_K\), where
\(K\) is odd, \(p\) is an odd prime, and \(p\nmid K\). Let \(L\)
be the least disputed size. For each \(0\le t\le a\), let
\(d_t:\mathbb Z_K\to\{-1,0,1\}\) indicate the signed size-\(L\)
disputes in the physical \(p\)-order state \(p^t\). If the functions
\(d_t\) are not all identical, their Fourier transforms have a common
zero at a nontrivial conductor of \(K\).
\end{lemma}

Here ``nontrivial'' means conductor greater than one. If \(Z\) denotes
the set of common zero conductors, the conclusion is
\(Z\setminus\{1\}\ne\varnothing\), rather than merely
\(Z\ne\varnothing\).

\begin{proof}
At a dyadic character of kernel order \(L\), divide the partial
profiles by \(L\), and denote their difference and sum in section
\(t\) by \(d_t\) and \(g_t\). The difference is exactly the stated
signed indicator: larger fibers contribute zero to this character,
smaller disputed fibers do not exist, and equal fibers cancel.
The functions \(g_t\) take values in \(\mathbb Z[1/2]\). For a section
vector \(u=(u_0,\ldots,u_a)\), put \(u_{a+1}=0\) and
\begin{equation}
T_k(u)=\sum_{i=0}^kp^i(u_i-u_{i+1})
      =\sum_{i=0}^k\varphi(p^i)u_i-p^ku_{k+1}.
\label{eq:primary-T}
\end{equation}
These are its Fourier rows on the \(p\)-axis, indexed by kernel
order \(p^k\).

Let \(j<a\) be the first change, so
\(d_0=\cdots=d_j\ne d_{j+1}\), and let \(Z\) be the set of
common zero conductors of all \(d_t\). At every residual conductor
outside \(Z\), invertibility of the \(p\)-axis transform gives some
nonzero row \(T_k(d)\). All rows below \(j\) vanish, hence \(k\ge j\).
Homometry forces the corresponding row \(T_k(g)\) to vanish. Moreover
\(T_k(g)\equiv T_j(g)\pmod{p^{j+1}}\). Thus the Fourier transform
of \(T_j(g)\) vanishes modulo \(p^{j+1}\) outside \(Z\).
Radial Fourier inversion on \(K\) uses integer Ramanujan sums and
division by \(K\). Since \(p\nmid 2K\), it is valid modulo
\(p^{j+1}\) in \(\mathbb Z[1/2]\). If \(Z\) were empty,
\(T_j(g)\) would be zero modulo \(p^{j+1}\); if \(Z\) consisted
only of the trivial conductor, \(T_j(g)\) would be constant modulo
\(p^{j+1}\).

We first record the argument when \(L\) is the only disputed size.
Intersect both sets with the subgroup having dyadic order \(L\).
Outside this subgroup the two sets agree, and shifts inside it
preserve every exterior coset. Subtracting the equal exterior
contributions proves that the restricted sets are homometric.
In this paragraph, \(g_t\) denotes the normalized sum for the
restricted pair at its trivial dyadic character; the functions \(d_t\)
are unchanged. All its disputes still have size \(L\). Pure \(p\)-axis marginal
balance now shows that every nonzero \(d_t\) has both signs and
zero total population. At a shell where a size-\(L\) dispute
disappears between \(j\) and \(j+1\), the sums through \(j\)
are one and the next sum is zero. At a shell where a dispute appears,
the preceding sums are two and the next sum is one. In either case
\(T_j(g)=p^j\). If \(d_{j+1}\ne0\), its two signs force both
restricted fibers at the residual origin to have size \(L\) through
state \(j+1\), giving \(T_j(g)=0\) there. If \(d_{j+1}=0\),
then \(d_0=d_j\ne0\), and its two signs force both fibers to be
empty at the largest residual-order shell in every section, again
giving \(T_j(g)=0\). These values are distinct modulo \(p^{j+1}\).
The common zero set must therefore contain a nontrivial conductor.

Return to the original profiles in the general case, allowing larger
disputed sizes.
Opposite signs cannot replace one another at a fixed residual shell.
At a disappearing size-\(L\) dispute the same argument still gives
\(T_j(g)=p^j\). At a newly appearing dispute the next sum is one.
Before it, the occupied fiber has size at least \(L\). If its size
exceeds \(L\), the other fiber is empty or has the same larger size,
and their partial sum is zero. Once it has size \(L\), both fibers
must have size \(L\), since \(d_0=\cdots=d_j\) is zero at this
shell. Thus the preceding sums are a sequence of zeros followed by
twos. Formula~\eqref{eq:primary-T} gives precisely the possibilities
\begin{equation}
p^j,\qquad -p^j,\qquad p^j-2p^b\quad(0\le b<j).
\label{eq:primary-change-residues}
\end{equation}
All are nonzero modulo \(p^{j+1}\). Consequently \(Z\) is nonempty.

Suppose for contradiction that \(Z\) contains only the trivial
conductor. Each \(d_t\) then has zero total population. If
\(d_0\ne0\), its two signs make both fibers empty at the largest
residual-order shell from the first section onward. This gives
\(T_j(g)=0\), contradicting the constant residue and
\eqref{eq:primary-change-residues}. We therefore have
\(d_0=\cdots=d_j=0\), while \(d_{j+1}\) has both signs.
At the residual origin both fibers are nonempty and equal, of size at
least \(L\), through state \(j+1\). Their partial sums are again
zeros followed by twos. Their value of \(T_j(g)\) is either zero
or \(-2p^c\) with \(0\le c\le j\).

Compare these origin residues with \eqref{eq:primary-change-residues}.
For \(b<j\), congruence of \(p^j-2p^b\) with \(-2p^c\)
first forces \(b=c\) by valuation and then would force
\(p^{j+1}\mid p^j\). Congruence of \(\pm p^j\) with
\(-2p^c\) requires \(c=j\) and \(\pm1\equiv-2\pmod p\).
The only possibility is \(p=3\) and the positive sign.
In this exceptional case, the origin fibers exceed \(L\) through
state \(j\) and first attain \(L\) at state \(j+1\); every newly
disputed shell has both fibers equal to \(L\) at every earlier state.

There can be no larger dispute in this exceptional case. At and after
state \(j+1\) the origin bounds all fibers by \(L\). If an earlier
section had a larger dispute at a residual shell \(s\), then all
preceding occupied fibers there would be larger than \(L\) and all
preceding partial sums would be zero. From that dispute to state
\(j\), its empty side stays empty, and the other side is either
larger than \(L\) or empty: a size-\(L\) dispute is excluded by
\(d_0=\cdots=d_j=0\), and smaller disputes are excluded by minimality.
Hence every partial sum through \(j\) at \(s\) is zero. At
\(j+1\) the sum is zero or one, yielding \(T_j(g)=0\) or
\(-p^j\). Neither is congruent to \(p^j\) modulo \(3^{j+1}\).
Thus only one disputed size occurs, and the one-size case already
proved gives the required contradiction.
\end{proof}

\begin{theorem}[An arbitrary odd prime power and three residual primes]
\label{thm:primary-power-four-primes}
Let \(p,q,r,s\) be pairwise distinct odd primes. For every
\(a\ge1\) and \(e\ge0\), the modulus \(2^ep^aqrs\) is rigid
among unions of subgroups.
\end{theorem}

\begin{proof}
Suppose a distinct homometric pair exists. Index the residual primes
\(q,r,s\) by \(1,2,3\), write \(u_i=q_i-1\), and label their
shells by subsets \(S\subseteq\{1,2,3\}\). Let
\(F_t(S)=h_A(t,S)-h_B(t,S)\), where the \(p\)-order is \(p^t\).
For fixed \(S\), its nonzero values have one sign: once a fiber is
empty it stays empty as \(t\) increases. Within any section \(t\),
opposite signs cannot occur at comparable residual shells.
Pure \(p\)-axis marginal balance gives
\begin{equation}
\sum_S\Bigl(\prod_{i\in S}u_i\Bigr)F_t(S)=0
\quad\hbox{for every }t.
\label{eq:primary-section-balance}
\end{equation}
Every disputed section therefore has both signs. Its bottom and top
shells cannot be disputed, since each is comparable with every shell.

Set
\[
W_S=\Bigl(\prod_{i\in S}u_i\Bigr)
       \sum_{t=0}^a\varphi(p^t)F_t(S).
\]
Because the sign is fixed along the primary axis, \(W_S=0\) if and
only if \(S\) is never disputed. Pure residual-axis marginals and
total population balance give
\(\sum_{S\ni i}W_S=0\) for each \(i\) and \(\sum_SW_S=0\).
Only singleton and doubleton shells can contribute. Summing the three
marginal equations and subtracting the total equation shows that the
sum over doubletons is zero; the sum over singletons is then zero too.
It follows that
\begin{equation}
W_{\{i\}}=W_{\{1,2,3\}\setminus\{i\}}.
\label{eq:primary-complementary-signs}
\end{equation}
Thus each singleton and its complementary doubleton are either both
inactive or have the same fixed sign.

In any section, every disputed singleton has the same sign as every
disputed doubleton: the two shells are either comparable or
complementary. If both ranks occurred, all disputes in the section
would have the same sign, contradicting
\eqref{eq:primary-section-balance}. Every disputed section lies
entirely in one rank and contains at most three disputed shells.

Let \(L\) be the least disputed size and let \(d_t\) indicate its
signed residual shells. These sections cannot all coincide. Otherwise
some size-\(L\) shell would occur in every section, while its
complementary shell must be disputed somewhere by
\eqref{eq:primary-complementary-signs}. That section would contain
both ranks. Lemma~\ref{lem:primary-common-zero} therefore supplies a
common nontrivial residual Fourier zero for all \(d_t\).
Write this character as \(\chi=\chi_1\otimes\chi_2\otimes\chi_3\),
with conductor \(c\mid q_1q_2q_3\). The lemma gives \(c>1\).
The local nonzero-shell sums are
\[
 x_i(c)=\sum_{z\in C_{q_i}\setminus\{0\}}\chi_i(z)
 =\begin{cases}
 q_i-1=u_i,&q_i\nmid c,\\
 -1,&q_i\mid c.
 \end{cases}
\]
For a nontrivial local character, the sum over the whole prime-order
group is zero; removing its value one at the origin leaves \(-1\).
Since \(c>1\), some \(q_i\mid c\), and hence some \(x_i=-1\).
In kernel coordinates this is \(c=\prod_iq_i^{1-h_i}\) with some
\(h_i=0\). The all-\(u_i\) evaluation occurs only at the trivial
character, whose kernel vector is \((1,1,1)\); the preceding lemma
explicitly excludes a common zero set consisting only of that character.
A physical shell \(S\) contributes \(\prod_{i\in S}x_i(c)\), which
is nonzero. Each section containing a size-\(L\) dispute must therefore
contain at least two of them.

Every such section must also contain a larger dispute. Indeed, suppose
all its disputes had size \(L\). Two opposite terms cannot have equal
population: the three shell populations within either rank are
pairwise distinct. Thus all three shells would be disputed, with both
signs. In rank one their polynomial is
\(\sum_i\epsilon_ix_i\), with \(\epsilon_i\in\{-1,1\}\).
At a nontrivial zero, parity forces exactly two coordinates to be
\(-1\). If those two signs differ, their contributions cancel and
the third is nonzero. If their signs agree, the third sign is opposite,
so all three evaluated contributions have the same sign. Either way
the polynomial cannot vanish. In rank two, divide by the nonzero
product \(x_1x_2x_3\); the polynomial becomes
\(\sum_i\epsilon_i/x_i\). With one coordinate equal to \(-1\),
the other two reciprocal magnitudes sum to at most
\(1/2+1/4<1\). With two such coordinates their signed sum is zero
or has magnitude two, and the remaining term has magnitude in
\((0,1/2]\). With three, the sum is an odd integer. This again
contradicts the common nontrivial zero.

Each section containing size \(L\) consequently has exactly two
size-\(L\) shells and one larger disputed shell. The size-\(L\)
shells have opposite signs. In rank one their polynomial is
\(\pm x_i\pm x_j\); in rank two it is that expression multiplied
by a common nonzero \(x_k\). A same-sign pair cannot vanish, since
\(x_i+x_j\ne0\) when each entry is \(-1\) or a positive even
integer.

All three complementary shell pairs are active. Their three fixed
signs include both signs, so exactly one sign is repeated. In every
section containing size \(L\), its two size-\(L\) shells have
opposite signs; the larger disputed shell therefore has the repeated
sign, independently of the section and its rank. Equation
\eqref{eq:primary-section-balance} forces the signed population of
the size-\(L\) shells in that section to have the opposite sign.
All nonzero such populations thus have the same sign. Their sum,
weighted by the positive numbers \(\varphi(p^t)\), cannot be zero.
This contradicts the exact-size population balance in
Lemma~\ref{lem:common-dyadic-fibers}, conclusion 3.
\end{proof}

\paragraph{Proof status.}
The preceding theorem is proved by the displayed arguments for every
prime and exponent in its statement. It uses no finite enumeration.
The hypothesis that the three residual primes have exponent one is
essential to the proof's six-shell argument; no conclusion for four
arbitrary prime-power axes is being inferred.
